% ============================================================
%  Does Labour Market Tightness Close the Gender Pay Gap?
%  A Panel Test of a Contested Premise
%  Author: Marti Soura Vamseekar
%  Overleaf-ready single file — compile with pdflatex + bibtex
%  Figures: pulled automatically from paper/exports/figures/ via
%           \graphicspath below. On Overleaf, upload this whole
%           paper/ folder (main.tex + exports/figures/) so the
%           relative path resolves.
% ============================================================

\documentclass[12pt,a4paper]{article}

% ---- Encoding & fonts ----------------------------------------
\usepackage[utf8]{inputenc}
\usepackage[T1]{fontenc}
\usepackage{lmodern}

% ---- Page layout ---------------------------------------------
\usepackage[a4paper,
            top=2.5cm, bottom=2.5cm,
            left=2.8cm, right=2.8cm]{geometry}

% ---- Line spacing (double for review) -----------------------
\usepackage{setspace}
\doublespacing

% ---- Math ----------------------------------------------------
\usepackage{amsmath,amssymb}

% ---- Tables --------------------------------------------------
\usepackage{booktabs}
\usepackage{array}
\usepackage{multirow}
\usepackage{threeparttable}   % table notes (Notes: ...)
\usepackage{adjustbox}        % rescale wide tables

% ---- Figures -------------------------------------------------
\usepackage{graphicx}
\graphicspath{{exports/figures/}{./}}
\usepackage{float}
\usepackage{caption}
\captionsetup{font=small, labelfont=bf, skip=6pt}

% ---- Typography ----------------------------------------------
\usepackage{microtype}

% ---- Hyperlinks ----------------------------------------------
\usepackage[colorlinks=true,
            linkcolor=blue!55!black,
            citecolor=blue!55!black,
            urlcolor=blue!55!black,
            pdftitle={Does Labour Market Tightness Close the Gender Pay Gap?},
            pdfauthor={Marti Soura Vamseekar}]{hyperref}

% ---- Bibliography --------------------------------------------
\usepackage[round,authoryear,sort]{natbib}

% ---- Misc ----------------------------------------------------
\usepackage{xcolor}
\usepackage{parskip}          % paragraph spacing instead of indent

% ---- Custom commands -----------------------------------------
\newcommand{\note}[1]{\par\smallskip\footnotesize\textit{Note:} #1}
\newcommand{\sym}[1]{\ensuremath{^{#1}}}   % significance stars inline

% ============================================================
%  TITLE & AUTHOR
% ============================================================
\title{\Large\textbf{Does Labour Market Tightness Close the Gender Pay Gap?\\[4pt]
       A Panel Test of a Contested Premise}}

\author{Marti Soura Vamseekar\\[6pt]
        \small Independent Researcher\\
        \small M.Sc.\ Data Science, University of Greenwich, London, UK\\[4pt]
        \small ORCID:\;\href{https://orcid.org/0009-0005-4884-1292}{0009-0005-4884-1292}\\
        \small Correspondence:\;\href{mailto:martisoura@gmail.com}{martisoura@gmail.com}}

\date{July 2026}

% ============================================================
\begin{document}
\maketitle
\thispagestyle{empty}

% ---- JEL + Keywords ------------------------------------------
\vspace{1em}
\noindent\textbf{JEL Codes:} J31, J71, J16, J21, C33, C55, K31

\smallskip
\noindent\textbf{Keywords:} Gender Pay Gap; Pay Transparency Directive; Labour
Market Tightness; Panel Fixed Effects; Sector Heterogeneity; Compliance
Exposure; Eurostat

\newpage
\setcounter{page}{1}

% ============================================================
%  ABSTRACT
% ============================================================
\begin{abstract}
\noindent
The EU Pay Transparency Directive (2023/970/EU), transposed into national law by
June 2026, is implicitly premised on competitive equalisation theory: tight
labour markets, by raising the cost of discrimination, should narrow gender pay
gaps.  We test this premise using a panel of 27 EU member states and Eurostat
labour-market data spanning 2019--2024.  A pooled cross-sectional specification
with year fixed effects reproduces a strong positive association between
employment rate and the gender pay gap ($\hat{\beta}=0.637$, SE\,=\,0.155,
$p<0.001$, $N=159$) --- the opposite sign from what competitive equalisation
predicts.  Once country fixed effects absorb time-invariant national
confounders, the coefficient collapses to $\hat{\beta}=0.113$ (SE\,=\,0.269,
$p=0.676$); a Hausman test confirms fixed effects as the preferred specification
($\chi^{2}=27.99$, $p<0.001$).  This null result survives an extensive battery of
checks: time-window splits, outlier exclusion, an alternative tightness proxy,
controlling for the male--female employment participation gap, a formal test of
whether the relationship is conditional on collective bargaining coverage, and an
event-study cross-check using the differential severity of the 2020 COVID-19
lockdown shock --- a source of variation with no evident channel back to national
gender-discrimination institutions.  We conclude the pooled correlation reflects
between-country heterogeneity, not a within-country causal channel, and that
labour market tightness is not a reliable signal for anticipating Directive
compliance risk.  Three findings remain robust and policy-relevant: gender pay
gaps differ sharply and consistently by sector (ANOVA $F=105.78$, $p<0.001$;
Finance 24.96\% mean gap versus Construction $-2.72$\%); Finance and Information
\& Communication sectors show 100\% exposure to the Directive's 5\%
joint-pay-assessment threshold across the sample period; and a classifier trained
only on 2019--2022 data and tested on 2023--2024 confirms that sector membership
substantially and significantly outperforms tightness as an out-of-sample
predictor of compliance exposure (test AUC 0.805 vs.\ 0.560, $p<0.0001$).  We
propose sector-based exposure mapping, now validated out of sample, as a more
actionable compliance-triage tool than labour-market tightness.
\end{abstract}

\newpage

% ============================================================
%  1. INTRODUCTION
% ============================================================
\section{Introduction}
\label{sec:intro}

Directive (EU) 2023/970, transposed into national law by 7~June~2026, requires
employers with 100 or more workers to report gender pay gaps by worker category
and to justify, through a joint pay assessment, any unadjusted gap of 5\% or
more.  Employer-level disclosures under the Directive begin in 2027 for the
largest firms.  The Directive's design carries an implicit empirical premise:
that labour markets left to their own competitive pressure are not closing gender
pay gaps on their own, and that regulatory intervention is therefore necessary.
This premise rests, in turn, on an older and contested claim in labour economics
--- that tight labour markets, in which employers compete harder for scarce
workers, should erode the scope for pay discrimination and narrow gaps, following
the logic of competitive equalisation \citep{Becker1957}.  If that mechanism
worked reliably, one observable implication would be a negative association
between labour market tightness and the gender pay gap across countries and over
time.  If it does not, tightness should show no reliable relationship with the
gap, or an unstable one.

This paper tests that observable implication directly, using a panel of 27 EU
member states and annual and quarterly Eurostat labour-market indicators covering
2019--2024.  We deliberately structure the empirical section to show the
argument unfold rather than presenting only a final specification.  A pooled,
cross-sectional test --- the kind of evidence most readily available and most
commonly cited in policy discussion --- shows a positive association between
employment rate and the gender pay gap, the opposite of what competitive
equalisation predicts.  We then show that this pooled association is a
between-country pattern: it does not survive the introduction of country and year
fixed effects, which absorb the influence of time-invariant national
characteristics such as industrial structure, wage-setting institutions, and
legal regimes.  The coefficient shrinks by roughly 82\% and loses all
statistical significance.  We test whether this null result is fragile or robust
across five independent challenges --- narrower time windows, the exclusion of
outlier countries, an alternative measure of tightness, and control for a leading
candidate selection mechanism (the male--female employment participation gap) ---
and find the null holds in every case.

This paper makes three contributions.  First, it provides the first
panel-identified test we are aware of of the labour-market-tightness premise
implicit in the Directive's regulatory rationale, showing that a widely citable
cross-sectional correlation does not survive the identification strategy
required to support a causal or even a reliable predictive claim; this null
proves robust not only to the checks in Section~\ref{subsec:robustness} but also
to a cross-check using a source of variation with no evident channel back to
national gender-discrimination institutions (the differential severity of the
2020 COVID-19 lockdown shock, Section~\ref{subsec:covid_event}) and to a formal
test of whether the relationship is conditional on collective bargaining
coverage (Section~\ref{subsec:bargaining}).  Second, it introduces a sector-based
compliance-exposure map against the Directive's own 5\% joint-pay-assessment
threshold, showing that exposure is concentrated by sector (Finance and
Information \& Communication at 100\% exposure across the sample) rather than by
labour-market tightness, and validates this mapping out of sample: a classifier
trained only on 2019--2022 data substantially and significantly outperforms a
tightness-based classifier when tested on 2023--2024 data it never saw during
fitting, turning a descriptive pattern into a demonstrated predictive one
(Section~\ref{subsec:oos_validation}).  Third, it confirms that sector-level
heterogeneity in gender pay gaps --- unlike the tightness relationship --- is
large, statistically robust, and stable across every specification tested,
offering a concrete empirical anchor for compliance prioritisation independent
of the tightness question.

The remainder of the paper proceeds as follows.  Section~\ref{sec:data}
describes the data and empirical strategy.  Section~\ref{sec:results} presents
results in the order the argument requires: descriptive statistics, the
pooled-to-fixed-effects contrast, robustness checks, an exogenous cross-check
using COVID-19 lockdown stringency, mechanism controls, a test for institutional
moderation by collective bargaining coverage, sector heterogeneity, compliance
exposure and its out-of-sample validation, and a descriptive risk-quadrant
classification.  Section~\ref{sec:discussion} discusses mechanisms and policy
implications.  Section~\ref{sec:conclusion} concludes.  Appendix~\ref{sec:appendix}
reports a diagnostic validation of the Hiring Pressure Index, and
Appendix~\ref{sec:appendixB} documents a pre-registered design for testing the
Directive's staggered transposition timing once post-2024 outcome data becomes
available.

% ============================================================
%  2. DATA AND METHODS
% ============================================================
\section{Data and Methods}
\label{sec:data}

\subsection{Data sources and panel construction}
\label{subsec:data_sources}

The primary dataset is a panel of 27 EU member states drawn from Eurostat's
labour force survey and structure-of-earnings-survey products, accessed via the
Eurostat JSON-stat API and processed through a versioned dbt pipeline
(\texttt{analytics/} in the accompanying repository).  Six core series are used:
employment rate (\texttt{lfsi\_emp\_a}), unemployment rate (\texttt{lfsi\_emp\_a}),
job vacancy rate (\texttt{jvs\_q\_nace2}), gender pay gap by NACE sector
(\texttt{earn\_gr\_gpgr2}), labour market flows (\texttt{lfsi\_long\_q}), and
labour market slack (\texttt{lfsi\_sla\_q}).  Employment rate, unemployment rate,
and gender pay gap are annual series; job vacancy rate, labour market slack, and
labour flow signals are quarterly and are averaged within each calendar year to
align with the annual series.

\medskip
\noindent\textbf{Data ceiling.} Eurostat's gender pay gap series
(\texttt{earn\_gr\_gpgr2}) was last updated 26~February~2026, and its most
recent published reference year is 2024.  All specifications reported below are
therefore restricted to 2019--2024; 2025 country-year rows are excluded rather
than imputed, since the Directive-relevant outcome variable is not yet available
for that year in the Eurostat source.  This is expected annual-series publication
lag rather than a data-collection failure, and we do not attempt to fill it.

\medskip
\noindent\textbf{Coverage.} All 27 EU member states have valid, non-missing
gender pay gap data from 2022 onward.  Greece's gender pay gap series begins in
2022 in the Eurostat source; 2019--2021 therefore have 26 of 27 countries with
valid data (Greece is the sole exclusion in those years).  The primary
fixed-effects panel has 159 country-year observations across 27 countries, close
to but short of a fully balanced $27 \times 6 = 162$ cell panel because of
Greece's shorter series.

\subsection{Composite indices}
\label{subsec:composites}

Two composite indices from the underlying WorkforceGuard AI analytics pipeline
are used in the compliance-exposure and risk-quadrant analyses
(Sections~\ref{subsec:exposure}--\ref{subsec:quadrant}): the Hiring Pressure Index
(HPI) and the Equity Risk Score (ERS), both computed in
\texttt{mart\_semantic\_metrics.sql} and reported per country as a
latest-available-period snapshot (not a year-varying panel signal).  HPI is
constructed from vacancy rate, unemployment rate, labour market slack, and labour
flow signals, with no gender pay gap input:

\begin{align*}
\text{HPI}_\text{raw} &= \text{vacancy\_rate} \times 11
  + \max(0,\;9 - \text{unemployment\_rate}) \times 4 \\
  &\quad + \max(0,\;12 - \text{labour\_slack\_rate}) \times 2.8
  + \text{flow\_to\_employment} \times 0.9
  + \text{flow\_to\_inactivity} \times 0.6
\end{align*}

ERS is a direct, near-linear transform of the gender pay gap itself:
$\text{ERS}_\text{raw} = \text{pay\_gap} \times 5.5$, both clamped to
$[0, 100]$.  Because ERS is mechanically derived from the same variable used to
define Directive-threshold exposure, any ERS--exposure relationship in
Section~\ref{subsec:exposure} is expected by construction and is reported only as
a consistency check, not as a substantive finding; HPI, which shares no input
with the gender pay gap, is the informative comparison.

We separately validate HPI's weighting scheme against a principal components
analysis of its six underlying input signals (Appendix~\ref{sec:appendix}).
This validation is weaker than initially hypothesised (Spearman
$\rho = 0.657$, $p = 0.156$, not significant at conventional thresholds on the
full 27-country panel) and is not claimed as a validated methodological
contribution in this paper; it is reported transparently in the appendix as a
diagnostic, not a result the main contributions depend on.

\subsection{Empirical strategy}
\label{subsec:strategy}

We estimate five specifications of the relationship between labour market
tightness (proxied primarily by employment rate) and the gender pay gap:

\begin{enumerate}
  \item Pooled OLS with year fixed effects, standard errors clustered by country.
  \item Two-way (country and year) fixed effects panel regression, standard
        errors clustered by country.  This is the paper's main specification.
  \item Two-way fixed effects with additional controls (unemployment rate, job
        vacancy rate).
  \item First-difference estimation, an alternative estimator robust to certain
        forms of unit-root behaviour.
  \item Two-way fixed effects using employment rate lagged one year, to address
        the possibility that a persistent pay gap reduces female labour force
        participation and mechanically depresses the employment rate (reverse
        causality).
\end{enumerate}

We test the appropriateness of fixed effects against random effects using a
Hausman test, and test the joint significance of the country fixed effects
against a pooled null using an F-test.  We then subject the main fixed-effects
specification to four independent robustness checks: splitting the sample into
2019--2022 and 2021--2024 sub-windows; excluding Italy and Spain as potential
outliers; substituting job vacancy rate for employment rate as the tightness
proxy; and, in a dedicated mechanism-controls analysis, adding the male--female
employment participation gap (a standard proxy for the
\citealt{OlivettiPetrongolo2008} selection channel) as a covariate in both the
pooled and fixed-effects specifications.

\subsection{Related literature}
\label{subsec:lit}

This paper sits at the intersection of four literatures: the measurement of
labour market tightness, the cross-country gender pay gap literature and its
selection problem, the empirical evaluation of pay transparency legislation, and
the methodology of composite indicator construction.

\medskip
\noindent\textbf{Labour market tightness.}  Tightness is formally defined in
search-and-matching theory as the vacancy-to-unemployment ratio
\citep{MortensenPissarides1994}, with the standard prediction that a high ratio
strengthens worker bargaining power and should erode the scope for wage
discrimination as employers compete for scarce labour.  This prediction has
motivated a substantial central-bank and policy literature on composite tightness
measurement, including the Federal Reserve Bank of New York's Hiring Pressure and
Wage (HPW) Tightness Index \citep{HeisePearceWeber2024}, the Federal Reserve
Bank of Kansas City's Labour Market Conditions Indicator
\citep{HakkioWillis2014}, and comparable composite indices from the Reserve Bank
of Australia \citep{BakerBall2018} and the Reserve Bank of New Zealand
\citep{ArmstrongKamberKaragedikli2016}.  At the EU level,
\citet{KissEtAl2022} analyse simultaneous slack and tightness across EU27 member
states post-COVID using vacancy rates, unemployment, and broader slack
indicators, estimating Beveridge curves to detect shifts in the
unemployment--vacancy relationship; their multi-signal approach directly motivates
this paper's use of a multi-signal Hiring Pressure Index rather than a single
tightness ratio.

\medskip
\noindent\textbf{The gender pay gap, selection, and sectoral variation.}  The
central theoretical tension for cross-country gender pay gap analysis was
formalised by \citet{OlivettiPetrongolo2008}, who show that gender wage gaps and
employment gaps are negatively correlated across countries: countries with lower
female employment rates exhibit smaller unadjusted pay gaps, because positive
selection into employment among only the most attached and highest-skill women
compresses the lower tail of the female wage distribution.
\citet{BlauKahn2017} provide the standard survey-level treatment of the gender
wage gap internationally.  \citet{ChristofidesPolycarpouVrachimis2013} extend
the selection analysis across the full earnings distribution using EU-SILC data,
documenting ``sticky floor'' and ``glass ceiling'' effects.
\citet{Hedija2017} provides sector-specific gender pay gap estimates across EU
member states using EU-SILC microdata, finding large cross-sector and
cross-country heterogeneity with Finance and insurance recording particularly
large residual gaps --- the closest precedent for this paper's sector
heterogeneity result (Section~\ref{subsec:sector}).
\citet{Bacovic2022} confirms persistent cross-country and sectoral heterogeneity
at a more aggregated level.

\medskip
\noindent\textbf{Pay transparency legislation: effects and conditions.}  The
empirical literature on pay transparency legislation has expanded substantially
since 2018.  \citet{BennedsenEtAl2022} exploit Denmark's 35-employee reporting
threshold using a difference-in-differences design and find treated firms reduced
their gender pay gap by approximately 2 percentage points relative to the
pre-legislation mean, driven primarily by slower male wage growth.
\citet{BlundellEtAl2025} extend this to the UK's mandatory gender pay gap
reporting regime, estimating a substantially larger effect.  A recurring finding
is that transparency legislation is markedly more effective where worker
representation structures exist to enforce it: \citet{VaccaroEtAl2024},
evaluating Germany's 2017 \textit{Entgelttransparenzgesetz}, find it reduced the
raw gender pay gap only in establishments with works councils and collective
bargaining agreements, with unexplained gaps in establishments lacking these
structures showing no improvement or even a slight increase.

\medskip
\noindent\textbf{Composite indicator methodology.}  The methodological basis for
constructing composite labour-market and social indicators is well established.
\citet{NardoEtAl2008} is the standard reference, identifying data-driven (PCA,
regression), normative, and equal-weighting strategies as the three principal
approaches, and noting that normative weights are often more defensible in
policy-oriented indices.  \citet{DecancqLugo2013} provide the standard overview
of weighting choices in multidimensional wellbeing indices, a framework directly
applicable to this paper's HPI and ERS.  This paper's Appendix~\ref{sec:appendix}
PCA validation of the HPI weighting scheme follows this literature's standard
diagnostic approach; consistent with \citeauthor{NardoEtAl2008}'s caution that
data-driven and normative weights need not agree, we find only a weak,
non-significant correlation.

\medskip
\noindent\textbf{What this paper adds.}  No published paper we are aware of
tests the labour-market-tightness premise implicit in the EU Pay Transparency
Directive using panel identification methods capable of distinguishing a
within-country relationship from a between-country correlation.
\citet{KissEtAl2022} analyse EU tightness without examining gender pay outcomes;
the gender pay gap literature \citep{OlivettiPetrongolo2008, Hedija2017,
BlauKahn2017} analyses pay gaps across countries without constructing or testing
against a composite tightness index.  This paper closes that specific gap, and
--- unusually for this literature --- reports a null result on the paper's own
central question after subjecting it to the full panel-identification and
robustness toolkit.

% ============================================================
%  3. RESULTS
% ============================================================
\section{Results}
\label{sec:results}

\subsection{Descriptive statistics}
\label{subsec:descriptive}

Table~\ref{tab:descriptive}, Panel~A reports variable-level descriptive
statistics across the 189 country-year cells in the full annual panel (27
countries, 7~years including 2025 for non-gender-pay-gap series).  Employment
rate averages 75.8\% (SD\,=\,5.1, range 58.3--83.6\%) with no missing
observations.  Unemployment rate averages 6.2\% (SD\,=\,2.7) with no missing
observations.  Gender pay gap, restricted to 2019--2024 as described in
Section~\ref{subsec:data_sources}, has 159 valid observations (15.9\% missing
across the full 189-cell frame, entirely attributable to the 2025 data ceiling
and Greece's shorter series) and averages 10.9\% (SD\,=\,5.6, range $-1.5$\% to
21.3\%).  Job vacancy rate has 3.7\% missing (182 of 189 cells) and averages
2.1\% (SD\,=\,1.2).

Table~\ref{tab:descriptive}, Panel~B reports country-level means sorted by mean
Hiring Pressure Index.  The Netherlands (HPI\,=\,100), Slovenia (HPI\,=\,99),
Germany (HPI\,=\,98), and Czechia (HPI\,=\,93) are the four tightest labour
markets in the sample by this measure; Romania (HPI\,=\,42), Greece (HPI\,=\,42),
and Finland (HPI\,=\,44) are among the loosest.

% ------- TABLE 1 -----------------------------------------------
\begin{table}[htbp]
\caption{Descriptive Statistics}
\label{tab:descriptive}
\begin{threeparttable}
\smallskip
\textit{Panel A: Variable-level summary statistics (full panel, $N=189$ country-year cells)}
\smallskip
\begin{adjustbox}{max width=\textwidth}
\begin{tabular}{lrrrrrrrrr}
\toprule
Variable & $N$ & Mean & SD & Min & P25 & Median & P75 & Max & \% Missing \\
\midrule
Employment rate & 189 & 75.80 & 5.14 & 58.30 & 72.80 & 77.00 & 79.70 & 83.60 & 0.00 \\
Unemployment rate & 189 & 6.15 & 2.68 & 2.00 & 4.40 & 5.80 & 7.10 & 17.90 & 0.00 \\
Gender pay gap & 159 & 10.92 & 5.56 & $-$1.50 & 7.00 & 11.40 & 15.60 & 21.30 & 15.87 \\
Job vacancy rate & 182 & 2.13 & 1.17 & 0.42 & 1.18 & 1.95 & 2.74 & 6.32 & 3.70 \\
Labour slack rate & 189 & 11.67 & 4.76 & 3.12 & 7.90 & 11.12 & 14.43 & 25.75 & 0.00 \\
\bottomrule
\end{tabular}
\end{adjustbox}

\bigskip
\textit{Panel B: Country-level means, sorted by Hiring Pressure Index (HPI)}
\smallskip
\begin{adjustbox}{max width=\textwidth}
\begin{tabular}{llrrrr}
\toprule
Code & Country & Mean Emp.\ Rate & Mean GPG (\%) & HPI & ERS \\
\midrule
NL & Netherlands  & 82.4 & 12.8 & 100 & 56 \\
SI & Slovenia     & 77.0 &  5.6 &  99 & 42 \\
DE & Germany      & 80.2 & 16.8 &  98 & 81 \\
CZ & Czechia      & 81.2 & 17.5 &  93 & 99 \\
CY & Cyprus       & 78.2 & 11.6 &  90 & 68 \\
AT & Austria      & 76.7 & 18.8 &  88 & 97 \\
HU & Hungary      & 79.6 & 16.3 &  88 & 90 \\
BE & Belgium      & 71.4 &  2.0 &  85 &  4 \\
PL & Poland       & 76.1 &  6.0 &  82 & 23 \\
BG & Bulgaria     & 75.2 & 12.3 &  81 & 63 \\
MT & Malta        & 79.6 &  6.9 &  80 & 27 \\
IE & Ireland      & 77.0 &  9.1 &  79 & 44 \\
LT & Lithuania    & 78.4 & 10.7 &  78 & 52 \\
LV & Latvia       & 77.1 & 16.9 &  76 & 71 \\
IT & Italy        & 64.8 &  5.3 &  71 & 34 \\
EE & Estonia      & 80.9 & 19.0 &  69 & 100\\
PT & Portugal     & 76.9 &  7.8 &  69 & 39 \\
HR & Croatia      & 70.1 &  7.9 &  68 & 32 \\
FR & France       & 73.8 & 13.7 &  67 & 67 \\
SE & Sweden       & 81.5 & 10.7 &  56 & 58 \\
DK & Denmark      & 79.3 & 13.5 &  55 & 74 \\
LU & Luxembourg   & 73.8 & $-$0.4 & 54 &  0 \\
ES & Spain        & 69.3 &  8.2 &  48 & 39 \\
SK & Slovakia     & 76.5 & 16.1 &  45 & 83 \\
FI & Finland      & 76.9 & 16.2 &  44 & 87 \\
EL & Greece       & 65.1 & 13.5 &  42 & 74 \\
RO & Romania      & 67.6 &  1.5 &  42 & 18 \\
\bottomrule
\end{tabular}
\end{adjustbox}
\begin{tablenotes}
\footnotesize
\item \textit{Notes:} Panel~A covers all 189 country-year cells (27~EU member
  states, 2019--2025).  Gender pay gap is restricted to 159 valid observations
  (2019--2024); missing cells reflect the 2025 Eurostat data ceiling and
  Greece's shorter series (2022--2024 only).  GPG\,=\,gender pay gap
  (percentage-point unadjusted gap, female earnings as \% below male).
  Panel~B reports means over 2019--2024; HPI and ERS are time-invariant
  latest-period scores (see Section~\ref{subsec:composites}).
  HPI range: 0--100 (higher = tighter market);
  ERS range: 0--100 (higher = larger pay gap).
\end{tablenotes}
\end{threeparttable}
\end{table}

\subsection{Pooled replication and fixed-effects null}
\label{subsec:pooled_fe}

Table~\ref{tab:regressions} reports the paper's central empirical result.
Model~1 (pooled OLS, year fixed effects) shows employment rate positively and
significantly associated with the gender pay gap: $\hat{\beta}=0.637$
(SE\,=\,0.155, $p<0.001$, $N=159$, $R^{2}=0.299$).  This replicates, on the
corrected and expanded 27-country panel, the positive cross-sectional association
widely noted in the policy literature.

Model~2, the paper's main specification, adds country and year fixed effects.
The coefficient falls to $\hat{\beta}=0.113$ (SE\,=\,0.269, $p=0.676$, $N=159$),
an attenuation of approximately 82\% relative to the pooled estimate, and the
association is not statistically distinguishable from zero at any conventional
threshold.  Within-country $R^{2}$ is $-0.096$, consistent with employment rate
explaining essentially none of the within-country variation in the gender pay gap
once between-country heterogeneity is absorbed.  A Hausman test strongly favours
fixed effects over random effects ($\chi^{2}=27.99$, $p<0.001$), and an F-test
confirms the country fixed effects are jointly significant against a pooled null
($F=83.80$, $p<0.001$) --- the attenuation in Model~2 is not an artefact of an
inappropriate estimator choice.

Model~3 (fixed effects with unemployment rate and job vacancy rate as additional
controls) gives $\hat{\beta}=0.209$ (SE\,=\,0.343, $p>0.05$, $N=153$),
consistent with Model~2.  Model~4 (first differences) gives
$\hat{\beta}=0.152$ (SE\,=\,0.112, $N=132$), and Model~5 (fixed effects with
employment rate lagged one year, addressing reverse causality) gives
$\hat{\beta}=0.016$ (SE\,=\,0.213, $N=132$) --- both null.  Every
fixed-effects-family specification in Table~\ref{tab:regressions} is
statistically indistinguishable from zero; only the pooled, cross-sectional
specification in Model~1 shows a significant association.

The contrast between Models~1 and~2 is the paper's central finding: a
relationship that appears strong and significant when identified only from
cross-country variation disappears entirely once within-country variation is
isolated.  This pattern is the signature of a confounded, between-country
association rather than a causal or even a reliably predictive within-country
relationship.  Figure~\ref{fig:scatter} illustrates the pooled cross-country
association.

% ------- TABLE 2 -----------------------------------------------
\begin{table}[htbp]
\caption{Labour Market Tightness and the Gender Pay Gap:\\
         Pooled OLS and Panel Fixed-Effects Estimates}
\label{tab:regressions}
\begin{threeparttable}
\begin{adjustbox}{max width=\textwidth}
\begin{tabular}{lcccccc}
\toprule
 & \multicolumn{5}{c}{\textit{Dependent variable: Gender pay gap (pp)}} \\
\cmidrule(l){2-6}
 & Model 1 & Model 2 & Model 3 & Model 4 & Model 5 \\
 & Pooled OLS & Two-way FE & FE + controls & First diff.\ & Lagged emp. \\
\midrule
Employment rate & $0.637^{***}$ & $0.113$ & $0.209$ & $0.152$ & $0.016$ \\
                & $(0.155)$     & $(0.269)$ & $(0.343)$ & $(0.112)$ & $(0.213)$ \\[4pt]
Country FE      & No  & Yes & Yes & No  & Yes \\
Year FE         & Yes & Yes & Yes & No  & Yes \\
Additional controls & No & No & Yes & No & No \\
\midrule
$N$               & 159 & 159 & 153 & 132 & 132 \\
$R^{2}$           & 0.299 & 0.007 & 0.038 & 0.018 & 0.000 \\
$R^{2}$ (within)  & ---  & $-$0.096 & $-$0.120 & --- & $-$0.009 \\
\bottomrule
\end{tabular}
\end{adjustbox}
\begin{tablenotes}
\footnotesize
\item \textit{Notes:} Standard errors clustered by country in parentheses.
  $^{***}p<0.001$, $^{**}p<0.01$, $^{*}p<0.05$.
  Model~1 is pooled OLS with year fixed effects.
  Model~2 (main specification) is two-way country-and-year fixed effects.
  Model~3 adds unemployment rate and job vacancy rate as controls.
  Model~4 is first-difference estimation.
  Model~5 uses employment rate lagged one year.
  Hausman test (RE vs.\ FE): $\chi^{2}=27.99$, $p<0.001$, favouring fixed effects.
  F-test for joint significance of country FEs in Model~2: $F=83.80$, $p<0.001$.
  Additional controls in Model~3: unemployment rate and job vacancy rate.
  Eurostat data ceiling: gender pay gap series ends at 2024 reference year
  (last updated 26~February~2026).
\end{tablenotes}
\end{threeparttable}
\end{table}

% ------- FIGURE 1 ----------------------------------------------
\begin{figure}[htbp]
  \centering
  \includegraphics[width=0.82\textwidth]{fig1_tightness_gpg_scatter}
  \caption{Labour Market Tightness and the Gender Pay Gap: Cross-Country Scatter.
    Each point represents one country-year observation (27~EU member states,
    2019--2024, $N=159$). The fitted line is from pooled OLS with year fixed
    effects (Model~1 in Table~\ref{tab:regressions},
    $\hat{\beta}=0.637^{***}$). The positive slope reflects a
    between-country pattern; once country fixed effects are introduced
    (Model~2), the coefficient collapses to 0.113 ($p=0.676$).}
  \label{fig:scatter}
\end{figure}

\subsection{Robustness checks}
\label{subsec:robustness}

Table~\ref{tab:robustness} subjects the Model~2 null result to four independent
robustness checks, all restricted to the same 2019--2024, 27-country panel.
Splitting the sample by time window, the coefficient is positive and insignificant
in 2019--2022 ($\hat{\beta}=0.300$, SE\,=\,0.326, $p=0.360$, $N=105$) and
negative and insignificant in 2021--2024 ($\hat{\beta}=-0.176$, SE\,=\,0.160,
$p=0.274$, $N=107$) --- the point estimate is unstable in sign across
sub-periods, and neither sub-period result is significant.  Excluding Italy and
Spain as potential outliers moves the coefficient only marginally
($\hat{\beta}=0.143$, SE\,=\,0.282, $p=0.614$, $N=147$), still null.
Substituting job vacancy rate for employment rate as the tightness proxy gives
$\hat{\beta}=-0.410$ (SE\,=\,0.270, $p=0.132$, $N=153$) --- the sign reverses
relative to the employment-rate specification, but this estimate is also not
statistically significant.  Every row in Table~\ref{tab:robustness} is null.
The overall robustness pattern is not ``one measure is significant and another is
not''; it is uniform non-significance across every specification and every proxy
we test, which is a stronger form of evidence against a stable tightness--pay-gap
relationship than a single null result would be on its own.

% ------- TABLE 5 -----------------------------------------------
\begin{table}[htbp]
\caption{Robustness Checks}
\label{tab:robustness}
\begin{threeparttable}
\begin{tabular}{llcccc}
\toprule
Check & Specification & $\hat{\beta}$ & SE & $p$-value & $N$ \\
\midrule
\multicolumn{6}{l}{\textit{A. Full sample benchmark (Model 2)}} \\[2pt]
 & Full sample, 2019--2024 & 0.113 & 0.269 & 0.676 & 159 \\[6pt]
\multicolumn{6}{l}{\textit{B. Time-window splits}} \\[2pt]
 & 2019--2022 & 0.300 & 0.326 & 0.360 & 105 \\
 & 2021--2024 & $-$0.176 & 0.160 & 0.274 & 107 \\[6pt]
\multicolumn{6}{l}{\textit{C. Outlier exclusion}} \\[2pt]
 & Excluding Italy and Spain & 0.143 & 0.282 & 0.614 & 147 \\[6pt]
\multicolumn{6}{l}{\textit{D. Alternative tightness proxy}} \\[2pt]
 & Employment rate (benchmark) & 0.113 & 0.269 & 0.676 & 159 \\
 & Job vacancy rate            & $-$0.410 & 0.270 & 0.132 & 153 \\
\bottomrule
\end{tabular}
\begin{tablenotes}
\footnotesize
\item \textit{Notes:} All specifications are two-way (country + year) fixed
  effects.  Dependent variable: gender pay gap (percentage points).  Standard
  errors clustered by country.  No specification reaches statistical
  significance at any conventional threshold.
\end{tablenotes}
\end{threeparttable}
\end{table}

\subsection{COVID-19 lockdown stringency as an exogenous cross-check}
\label{subsec:covid_event}

The checks in Section~\ref{subsec:robustness} vary the sample, the outlier
set, and the tightness proxy, but all rely on ordinary year-to-year movement
in employment rate --- variation that could in principle still be correlated
with unobserved, time-varying national factors.  As a cross-check using
variation plausibly unrelated to a country's underlying discrimination
institutions, we exploit the 2020 COVID-19 lockdown shock: countries differed
sharply in lockdown stringency for reasons largely disconnected from their
gender-pay-gap institutions (public-health capacity, population density,
political response), yet the shock produced large, cross-country-variable
disruptions to employment during exactly the years this panel covers.

We construct $\text{Stringency2020}_i$, each country's mean 2020 daily value
on the Oxford COVID-19 Government Response Tracker's Stringency Index
\citep{HaleEtAl2021}, and estimate an event-study specification interacting
this time-invariant shock intensity with year dummies:

\begin{equation}
\text{GPG}_{it} = \sum_{k=2020}^{2024} \theta_k \left(\text{Stringency2020}_i
  \times \mathbb{1}[t=k]\right) + \gamma_i + \delta_t + \varepsilon_{it},
\label{eq:eventstudy}
\end{equation}

\noindent with country and year fixed effects as in Model~2, standard errors
clustered by country, and 2019 as the omitted reference year.  If tightness
genuinely drove the gender pay gap, countries hit harder by the 2020 shock
should show a divergent path in $\theta_k$ relative to less-affected
countries.

Table~\ref{tab:eventstudy} and Figure~\ref{fig:covid_event} report the
result: no $\theta_k$ is statistically distinguishable from zero in any year,
2020--2024.  The point estimates drift from near zero in 2020
($\hat\theta=-0.012$, $p=0.741$) to modestly positive by 2024
($\hat\theta=0.085$, $p=0.190$), but every confidence interval spans zero
throughout, and the widening intervals in the later years reflect
accumulating imprecision rather than a strengthening relationship.  We note
one limitation of this design explicitly: because the panel begins in 2019,
there is only a single pre-treatment year, so we cannot formally test for
pre-existing divergent trends between high- and low-stringency countries
before the shock.  The design's identifying assumption --- that stringency
is unrelated to gender-pay-gap institutions conditional on country and year
effects --- is therefore plausible but not independently verifiable within
this sample.

This result is a stronger form of null than the observational checks in
Section~\ref{subsec:robustness}, because it uses variation --- an
exogenous public-health shock --- that has no obvious channel back to a
country's wage-setting institutions, rather than ordinary employment-rate
fluctuation that could in principle share unobserved drivers with the gender
pay gap.\footnote{We also considered a two-stage least-squares specification
instrumenting employment rate with $\text{Stringency2020}_i \times
\text{Post}_t$.  This is not informative in a two-way fixed-effects design:
$\text{Post}_t$ is a deterministic function of the year fixed effects already
in the model, leaving almost no independent variation to identify the
instrument (partial $F=0.627$, $p=0.430$).  We therefore rely on the
event-study estimates above, which do not require this exclusion
restriction, rather than reporting the 2SLS estimate as a robustness
confirmation.}

% ------- TABLE: COVID EVENT STUDY --------------------------------
\begin{table}[htbp]
\caption{Event Study: COVID-19 Lockdown Stringency and the Gender Pay Gap}
\label{tab:eventstudy}
\begin{threeparttable}
\begin{tabular}{lccccc}
\toprule
Year & $\hat\theta_k$ & SE & $p$-value & 95\% CI (low) & 95\% CI (high) \\
\midrule
2019 (ref.) & 0.000  & --    & --    & --      & --    \\
2020        & $-$0.012 & 0.035 & 0.741 & $-$0.080 & 0.057 \\
2021        & 0.070  & 0.047 & 0.140 & $-$0.023 & 0.162 \\
2022        & 0.026  & 0.036 & 0.475 & $-$0.045 & 0.096 \\
2023        & 0.046  & 0.061 & 0.454 & $-$0.074 & 0.167 \\
2024        & 0.085  & 0.065 & 0.190 & $-$0.042 & 0.212 \\
\bottomrule
\end{tabular}
\begin{tablenotes}
\footnotesize
\item \textit{Notes:} Dependent variable: gender pay gap (percentage points).
  Two-way (country + year) fixed effects, standard errors clustered by
  country.  $\text{Stringency2020}_i$ is each country's mean 2020 Oxford
  COVID-19 Government Response Tracker Stringency Index value
  \citep{HaleEtAl2021}, interacted with year dummies; 2019 is the omitted
  reference year.  No coefficient is statistically significant at
  conventional thresholds.
\end{tablenotes}
\end{threeparttable}
\end{table}

% ------- FIGURE: COVID EVENT STUDY --------------------------------
\begin{figure}[htbp]
  \centering
  \includegraphics[width=0.82\textwidth]{fig_covid_event_study}
  \caption{Event Study: COVID-19 Lockdown Stringency and the Gender Pay Gap.
    Coefficients on $\text{Stringency2020}_i \times \mathbb{1}[t=k]$ from
    Equation~\ref{eq:eventstudy}, with 95\% confidence intervals (clustered by
    country).  2019 is the omitted reference year.  No year shows a
    statistically significant deviation from zero.}
  \label{fig:covid_event}
\end{figure}

\subsection{Mechanism controls}
\label{subsec:mechanism}

A natural question is whether the pooled association in Model~1 reflects a
genuine but confounded selection mechanism: countries with higher overall
employment rates might also have systematically different male--female labour
force participation gaps, which could affect the measured gender pay gap through
workforce composition rather than through a labour-demand channel
\citep{OlivettiPetrongolo2008}.  We test this directly by constructing a gender
employment gap variable (male employment rate minus female employment rate, both
age 20--64, Eurostat \texttt{lfsi\_emp\_a}) and adding it as a control in both
the pooled and fixed-effects specifications.  Employment rate (total) and female
employment rate are correlated at $r=0.95$ in this panel and cannot be included
together as separate regressors without severe multicollinearity; the gender
employment gap avoids this problem by construction.

Table~\ref{tab:mechanism}, Panel~A shows that adding the gender employment gap
to the pooled specification attenuates the employment-rate coefficient only
modestly, from $\hat{\beta}=0.637$ to $\hat{\beta}=0.578$ --- a 9.3\% reduction
--- and the coefficient remains highly significant ($p=0.003$).  Adding female
part-time share as a further control produces negligible additional change
($\hat{\beta}=0.643$ without the gender employment gap, $\hat{\beta}=0.584$
with both controls).  The gender employment participation gap does not
substantially explain the pooled cross-country association.

Table~\ref{tab:mechanism}, Panel~B repeats the exercise under country and year
fixed effects.  The baseline null result ($\hat{\beta}=0.113$, $p=0.676$) is
essentially unchanged when the gender employment gap is added
($\hat{\beta}=0.213$, SE\,=\,0.236, $p=0.368$, $N=159$): the point estimate
moves but remains well within the range of statistical noise, and the
specification remains clearly insignificant.  The gender employment gap alone,
without employment rate, is also insignificant within-country
($\hat{\beta}=0.140$, SE\,=\,0.206, $p=0.496$).

We report this as a negative result on its own terms: the leading candidate
selection mechanism we test does not explain the pooled correlation, and the
within-country null persists whether or not this mechanism is controlled for.
This does not rule out other unobserved between-country confounders ---
industrial structure, collective bargaining coverage, or childcare policy regimes
are plausible candidates we do not test directly --- but it forecloses one
specific, testable selection story.

% ------- TABLE 7 -----------------------------------------------
\begin{table}[htbp]
\caption{Mechanism Controls: Gender Employment Participation Gap}
\label{tab:mechanism}
\begin{threeparttable}
\smallskip
\textit{Panel A: Pooled OLS with year fixed effects}
\smallskip
\begin{tabular}{lcccc}
\toprule
Specification & $\hat{\beta}_{\text{emp.rate}}$ & SE & $p$-value & $N$ \\
\midrule
(a) Employment rate only [baseline]
  & $0.637^{***}$ & 0.155 & $<$0.001 & 159 \\
(b) + gender employment gap
  & $0.578^{**}$  & 0.196 & 0.003    & 159 \\
(c) + female part-time share
  & $0.643^{***}$ & 0.163 & $<$0.001 & 159 \\
(d) + gender employment gap + female part-time share
  & $0.584^{**}$  & 0.208 & 0.005    & 159 \\
\bottomrule
\end{tabular}

\bigskip
\textit{Panel B: Two-way (country + year) fixed effects}
\smallskip
\begin{tabular}{lcccc}
\toprule
Specification & $\hat{\beta}_{\text{emp.rate}}$ & SE & $p$-value & $N$ \\
\midrule
(e) Employment rate only [existing null]
  & $0.113$  & 0.269 & 0.676 & 159 \\
(f) + gender employment gap
  & $0.213$  & 0.236 & 0.368 & 159 \\
(g) Gender employment gap only (no employment rate)
  & $0.140$  & 0.206 & 0.496 & 159 \\
\bottomrule
\end{tabular}
\begin{tablenotes}
\footnotesize
\item \textit{Notes:} Dependent variable: gender pay gap (percentage points).
  Standard errors clustered by country.
  $^{***}p<0.001$, $^{**}p<0.01$, $^{*}p<0.05$.
  Gender employment gap\,=\,male minus female employment rate (ages 20--64,
  Eurostat \texttt{lfsi\_emp\_a}).  Female part-time share from Eurostat
  labour force survey.  Panel~A: pooled OLS; Panel~B: two-way FE.
  The 9.3\% attenuation from spec.~(a) to spec.~(b) in Panel~A is not
  statistically meaningful; the Olivetti--Petrongolo selection channel does
  not explain the pooled cross-country association.
\end{tablenotes}
\end{threeparttable}
\end{table}

\subsection{Does bargaining power moderate the tightness--gap relationship?}
\label{subsec:bargaining}

Section~\ref{subsec:mechanism} names collective bargaining coverage as a
plausible unobserved confound behind the pooled cross-country association,
alongside industrial structure and childcare policy regimes, but does not
test it directly.  There is also a specific, literature-motivated reason to
expect tightness to matter conditionally rather than not at all:
\citet{VaccaroEtAl2024} find that Germany's 2017 pay-transparency law reduced
the gender pay gap only in establishments with works councils and collective
bargaining coverage, with no improvement in establishments lacking these
structures.  If institutional voice is similarly necessary for a
competitive-equalisation channel to operate, tightness might compress the
gender pay gap where bargaining coverage is high while showing no effect
where it is low --- a pattern the pooled null in Section~\ref{subsec:pooled_fe}
cannot distinguish from a true absence of any tightness effect.

We construct $\text{BargCov}_i$, each country's adjusted collective
bargaining coverage rate from the OECD/AIAS ICTWSS database, version~2.0
\citep{OECDAIAS2025}, using the latest available pre-2019 value as a
time-invariant institutional baseline.\footnote{For four countries (Austria,
Belgium, France, and Italy), the OECD/AIAS database's own methodology notes
that the reported adjusted coverage rate is modelled rather than directly
measured from administrative or survey sources; we retain these countries for
comparability with the rest of the panel but flag this as a data-quality
caveat.}  Because $\text{BargCov}_i$ is time-invariant, it is absorbed by the
country fixed effects on its own; it is identified only through its
interaction with the within-country variation in employment rate:

\begin{equation}
\text{GPG}_{it} = \beta_1 \,\text{EmpRate}_{it}
  + \beta_2 \left(\text{EmpRate}_{it} \times \text{BargCov}_i\right)
  + \gamma_i + \delta_t + \varepsilon_{it}.
\label{eq:bargaining}
\end{equation}

Table~\ref{tab:bargaining}, Panel~A reports our pre-specified test.  The
interaction term is small and statistically indistinguishable from zero
($\hat\beta_2=-0.0029$, SE\,=\,0.0026, $p=0.280$, $N=159$), and the
employment-rate main effect remains null as in Model~2
($\hat\beta_1=0.228$, $p=0.388$).  On this test, the tightness--gap null does
not appear to be institutionally gated: it holds regardless of a country's
baseline bargaining coverage.

As a secondary, more exploratory cut, Panel~B replaces the continuous
interaction with a binary above/below-median split on $\text{BargCov}_i$ and
formally tests whether the implied slopes differ.  Here the difference is
statistically significant ($p=0.023$): the implied slope is negative among
above-median-coverage countries ($-0.112$) and positive, though
insignificant on its own, among below-median-coverage countries ($0.309$,
$p=0.233$) --- directionally consistent with competitive equalisation
operating only where bargaining coverage is high, the pattern
\citet{VaccaroEtAl2024} document for transparency legislation.  We report
this cut honestly rather than suppress it, but we do not treat it as
confirming institutional gating: it disagrees with our pre-specified
continuous test, binary median-splits of continuous variables are known to
be sensitive to the exact cutpoint chosen, and each half of the sample
contains only 13--14 country clusters, which is thin for cluster-robust
inference.  We flag the binary-split pattern as a candidate for future
research using a longer panel or firm-level bargaining data, in the same
spirit as our treatment of the exploratory ICT-sector result in
Section~\ref{subsec:sector}, rather than as a second confirmed finding
alongside the null in Panel~A.

% ------- TABLE: BARGAINING COVERAGE INTERACTION ------------------
\begin{table}[htbp]
\caption{Does Collective Bargaining Coverage Moderate the Tightness--Gap
         Relationship?}
\label{tab:bargaining}
\begin{threeparttable}
\smallskip
\textit{Panel A: Continuous interaction (primary, pre-specified test)}
\smallskip
\begin{tabular}{lccc}
\toprule
 & Coefficient & SE & $p$-value \\
\midrule
Employment rate                          & 0.228    & 0.264  & 0.388 \\
Employment rate $\times$ BargCov         & $-$0.0029 & 0.0026 & 0.280 \\
\midrule
$N$ & \multicolumn{3}{r}{159 (27 countries)} \\
\bottomrule
\end{tabular}

\bigskip
\textit{Panel B: Binary median-split, formal test of slope difference (secondary)}
\smallskip
\begin{tabular}{lccc}
\toprule
 & Coefficient & SE & $p$-value \\
\midrule
Employment rate, below-median BargCov         & 0.309    & 0.258 & 0.233 \\
Difference in slopes (above $-$ below median) & $-$0.421 & 0.182 & 0.023 \\
Implied slope, above-median BargCov           & $-$0.112 & ---   & --- \\
\midrule
$N$ & \multicolumn{3}{r}{159 (27 countries); median BargCov\,=\,52.4} \\
\bottomrule
\end{tabular}
\begin{tablenotes}
\footnotesize
\item \textit{Notes:} Dependent variable: gender pay gap (percentage points).
  Two-way (country + year) fixed effects, standard errors clustered by
  country.  BargCov: adjusted collective bargaining coverage rate,
  OECD/AIAS ICTWSS v2.0 \citep{OECDAIAS2025}, latest pre-2019 value,
  time-invariant.  Panel~A is the pre-specified test; Panel~B is a secondary
  robustness cut with 13--14 country clusters per half and should be
  interpreted cautiously (see text).
\end{tablenotes}
\end{threeparttable}
\end{table}

\subsection{Sector heterogeneity}
\label{subsec:sector}

In sharp contrast to the fragile tightness relationship, gender pay gaps differ
substantially and robustly across sectors.  Restricting to single-letter NACE
Rev.~2 sections to avoid double-counting sectors against their own aggregate
rollups, an ANOVA across 19 disjoint sectors rejects the null of equal means
decisively ($F=105.78$, $p<0.001$).  Financial and insurance activities has the
highest mean gap in the sample at 24.96\% (SD\,=\,6.62, $N=154$), and
Construction has the lowest at $-2.72$\% (SD\,=\,11.51, $N=154$) --- a spread
of nearly 28 percentage points between the two extremes.  Every pairwise
comparison of Finance against every other sector remains significant after
Bonferroni correction for multiple comparisons (full pairwise results in
Appendix~\ref{sec:appendixC}).  This sector effect is present
in every year of the sample and is not sensitive to the specification choices
that eliminate the tightness relationship.  Figure~\ref{fig:sectorbars}
illustrates the sector distribution.

The country-level null reported in Section~\ref{subsec:pooled_fe}
($\hat{\beta}=0.113$) does not by itself rule out a tightness relationship
operating differently within specific sectors, so we re-estimate the main
two-way fixed-effects specification separately for Finance and for Information
and Communication, the two highest-exposure sectors identified below
(Table~\ref{tab:sectorfe}).  The Finance-sector coefficient is small and
statistically insignificant ($\hat{\beta}=-0.047$, SE\,=\,0.343, $p=0.891$,
$N=154$), consistent with the country-level null.  The ICT-sector coefficient,
by contrast, is larger and nominally significant ($\hat{\beta}=0.832$,
SE\,=\,0.391, $p=0.035$, $N=154$).  We report this result cautiously rather
than as a second confirmed finding: it is one of only two sector-specific
specifications we test, so it has not been subjected to a multiple-comparisons
correction, and NACE sector composition differs enough across countries that a
single-sector fixed-effects estimate is more susceptible to composition effects
than the pooled country-level result.  We do not treat this result as overturning
the paper's main country-level null claim; it is noted here as a candidate for
future sector-specific investigation rather than as an established finding.

% ------- TABLE 3 -----------------------------------------------
\begin{table}[htbp]
\caption{Sector Heterogeneity in Gender Pay Gaps}
\label{tab:sector}
\begin{threeparttable}
\begin{adjustbox}{max width=\textwidth}
\begin{tabular}{p{8.5cm}crrrrr}
\toprule
Sector & Code & Mean GPG & SD & Min & Max & $N$ \\
\midrule
Financial and insurance activities         & K & 24.96 & 6.62 & 8.4 & 40.3 & 154 \\
Information and communication              & J & 19.46 & 6.66 & 6.5 & 30.8 & 154 \\
Wholesale and retail trade                 & G & 18.23 & 5.73 & 4.2 & 32.9 & 154 \\
Human health and social work activities    & Q & 17.07 & 8.57 & $-$5.4 & 33.0 & 154 \\
Manufacturing                              & C & 16.75 & 7.08 & 0.9 & 30.9 & 154 \\
Arts, entertainment and recreation         & R & 16.55 & 13.58 & $-$8.4 & 67.3 & 154 \\
Professional, scientific and technical act.\ & M & 16.31 & 7.12 & $-$5.4 & 28.1 & 154 \\
Other service activities                   & S & 12.84 & 8.34 & $-$10.1 & 31.3 & 154 \\
Real estate activities                     & L & 11.47 & 8.11 & $-$21.3 & 36.7 & 154 \\
Education                                  & P &  9.95 & 5.73 & $-$6.4 & 23.5 & 154 \\
Accommodation and food service activities  & I &  9.22 & 4.29 & 1.5 & 28.0 & 154 \\
Electricity, gas, steam and air supply     & D &  8.48 & 7.43 & $-$9.1 & 23.5 & 142 \\
Public administration and defence          & O &  8.07 & 5.28 & $-$4.1 & 20.2 & 126 \\
Administrative and support service act.\   & N &  7.43 & 9.99 & $-$30.3 & 24.7 & 154 \\
Mining and quarrying                       & B &  6.94 & 10.48 & $-$18.5 & 39.2 & 148 \\
Transportation and storage                 & H &  3.79 & 8.68 & $-$21.0 & 32.3 & 154 \\
Water supply; sewerage \& waste management & E & $-$0.28 & 12.28 & $-$54.8 & 20.9 & 152 \\
Construction                               & F & $-$2.72 & 11.51 & $-$31.2 & 28.2 & 154 \\
Agriculture, forestry and fishing          & A & \multicolumn{5}{c}{--- (no usable observations)} \\
\bottomrule
\end{tabular}
\end{adjustbox}
\begin{tablenotes}
\footnotesize
\item \textit{Notes:} GPG\,=\,unadjusted gender pay gap (percentage points,
  female earnings as \% below male).  Single-letter NACE Rev.~2 sections only;
  Agriculture (A) has no usable observations in the Eurostat source.
  ANOVA across 18 sectors (excluding A): $F=105.78$, $p<0.001$.
  Every pairwise comparison of Finance (K) against all other sectors is
  significant after Bonferroni correction.
  Sample: 27~EU member states, 2019--2024.
\end{tablenotes}
\end{threeparttable}
\end{table}

% ------- FIGURE 3 ----------------------------------------------
\begin{figure}[htbp]
  \centering
  \includegraphics[width=0.88\textwidth]{fig3_sector_gpg_bars}
  \caption{Gender Pay Gaps by Sector. Mean unadjusted gender pay gap
    (percentage points) by NACE Rev.~2 single-letter section across 27~EU
    member states, 2019--2024. Error bars show $\pm$1~standard deviation.
    ANOVA $F=105.78$, $p<0.001$.}
  \label{fig:sectorbars}
\end{figure}

% ------- TABLE 3c ----------------------------------------------
\begin{table}[htbp]
\caption{Sector-Specific Two-Way Fixed-Effects Estimates}
\label{tab:sectorfe}
\begin{threeparttable}
\begin{tabular}{lccccl}
\toprule
Sector & $\hat{\beta}$ & SE & $p$-value & $N$ & Country-level $\hat{\beta}$ \\
\midrule
Financial and insurance activities & $-$0.047 & 0.343 & 0.891 & 154 & 0.113 \\
Information and communication      & $0.832^{*}$ & 0.391 & 0.035 & 154 & 0.113 \\
\bottomrule
\end{tabular}
\begin{tablenotes}
\footnotesize
\item \textit{Notes:} Both specifications are two-way (country + year) fixed
  effects, estimated on the sector-specific subsample. Dependent variable:
  gender pay gap for the named sector.  Standard errors clustered by country.
  $^{*}p<0.05$.  The country-level $\hat{\beta}=0.113$ (Model~2,
  Table~\ref{tab:regressions}) is shown for comparison.  The ICT estimate
  should be interpreted cautiously: it is one of two sector-specific
  specifications tested and has not been subjected to a multiple-comparisons
  correction.
\end{tablenotes}
\end{threeparttable}
\end{table}

\subsection{Compliance exposure under the Directive's 5\% threshold}
\label{subsec:exposure}

Directive (EU) 2023/970 requires a joint pay assessment when an unadjusted
gender pay gap of 5\% or more is not objectively justified.  We compute, for
every country-sector-year cell in the 2019--2024 panel (excluding 2025 per
Section~\ref{subsec:data_sources}), whether the observed gender pay gap meets or
exceeds this threshold, and examine how exposure varies by sector and by each
country's Hiring Pressure Index classification.

Of 4,863 country-sector-year cells in the primary panel, 3,173 (65.2\%) have
usable, non-missing gender pay gap data.  Restricting to Hiring Pressure Index
quartiles (a time-invariant, latest-period country classification independent of
the gender pay gap itself; see Section~\ref{subsec:composites}), exposure rises
from 75.9\% in the lowest-tightness quartile to 86.0\% in the highest-tightness
quartile (Q2:\,79.8\%, Q3:\,76.9\%).  This is the opposite of what a ``tight
markets self-correct'' story would predict: the tightest labour markets in the
sample do not show lower exposure to the Directive's threshold than the loosest
ones; if anything, exposure is somewhat higher in the tightest quartile.  We
deliberately lead with the HPI comparison rather than the parallel Equity Risk
Score comparison, since ERS is constructed directly from the gender pay gap
itself and an ERS--exposure relationship is close to tautological; the HPI
result, built from signals with no gender pay gap input, is the informative test.

Sector-level exposure (Table~\ref{tab:sectorexposure}, single-letter NACE
sections only) is far more discriminating than the tightness classification.
Financial and insurance activities and Information and Communication both show
100\% exposure --- every observed country-year cell in these two sectors meets or
exceeds the 5\% threshold across the sample period.  Wholesale and retail trade
follows at 98.1\%, and Professional, scientific and technical activities at
94.2\%.  At the country level (Table~\ref{tab:countryexposure}), Austria shows
100\% exposure across its 114 usable cells, followed by Estonia (94.4\%),
Czechia (93.7\%), and Latvia (92.9\%).

Sector, not labour-market tightness, is the variable that discriminates sharply
and consistently between high- and low-exposure cells in this data.  We propose
sector-based exposure mapping --- flagging Finance and ICT specifically --- as a
more actionable compliance-triage heuristic for regulators and employers than any
classification built on labour-market tightness.  Figure~\ref{fig:exposure}
illustrates sector-level exposure.

% ------- TABLE 6 -----------------------------------------------
\begin{table}[htbp]
\caption{Compliance Exposure by Hiring Pressure Index and Equity Risk Score
         Quartile}
\label{tab:exposure}
\begin{threeparttable}
\begin{adjustbox}{max width=\textwidth}
\begin{tabular}{llrrr}
\toprule
Index & Quartile & $N$ usable & $N \geq 5\%$ & Exposure (\%) \\
\midrule
\multicolumn{5}{l}{\textit{HPI (full panel, time-invariant classification)}} \\[2pt]
 & Q1 (lowest)  & 801 & 608 & 75.9 \\
 & Q2           & 868 & 693 & 79.8 \\
 & Q3           & 874 & 672 & 76.9 \\
 & Q4 (highest) & 630 & 542 & 86.0 \\[4pt]
\multicolumn{5}{l}{\textit{ERS (full panel, time-invariant classification)}} \\[2pt]
 & Q1 (lowest)  & 838 & 551 & 65.8 \\
 & Q2           & 780 & 604 & 77.4 \\
 & Q3           & 813 & 692 & 85.1 \\
 & Q4 (highest) & 742 & 668 & 90.0 \\[4pt]
\multicolumn{5}{l}{\textit{HPI (2024 only, robustness)}} \\[2pt]
 & Q1 (lowest)  & 143 & 108 & 75.5 \\
 & Q2           & 143 & 111 & 77.6 \\
 & Q3           & 142 & 103 & 72.5 \\
 & Q4 (highest) & 105 &  91 & 86.7 \\[4pt]
\multicolumn{5}{l}{\textit{ERS (2024 only, robustness)}} \\[2pt]
 & Q1 (lowest)  & 139 &  86 & 61.9 \\
 & Q2           & 127 &  99 & 78.0 \\
 & Q3           & 145 & 119 & 82.1 \\
 & Q4 (highest) & 122 & 109 & 89.3 \\
\bottomrule
\end{tabular}
\end{adjustbox}
\begin{tablenotes}
\footnotesize
\item \textit{Notes:} Exposure is defined as gender pay gap $\geq 5\%$
  (Directive joint-pay-assessment threshold).  HPI and ERS quartiles are
  time-invariant, latest-period country classifications.  ERS is mechanically
  derived from the gender pay gap ($\text{ERS}_\text{raw}=\text{pay\_gap}\times
  5.5$); the ERS--exposure relationship is therefore largely tautological and
  shown only for completeness.  HPI, which shares no input with the gender pay
  gap, is the informative comparison.
  Eurostat annual gender pay gap statistics (\texttt{earn\_gr\_gpgr2}) were last
  updated 26~February~2026; the latest available reference year is 2024.
\end{tablenotes}
\end{threeparttable}
\end{table}

% ------- TABLE 6b ----------------------------------------------
\begin{table}[htbp]
\caption{Sector-Level Compliance Exposure}
\label{tab:sectorexposure}
\begin{threeparttable}
\begin{adjustbox}{max width=\textwidth}
\begin{tabular}{p{8.5cm}rrrr}
\toprule
Sector & $N$ usable & $N \geq 5\%$ & Mean GPG (\%) & Exposure (\%) \\
\midrule
Financial and insurance activities                & 154 & 154 & 24.96 & \textbf{100.0} \\
Information and communication                     & 154 & 154 & 19.46 & \textbf{100.0} \\
Wholesale and retail trade                        & 154 & 151 & 18.23 & 98.1 \\
Professional, scientific and technical activities & 154 & 145 & 16.31 & 94.2 \\
Manufacturing                                     & 154 & 141 & 16.75 & 91.6 \\
Human health and social work activities           & 154 & 140 & 17.07 & 90.9 \\
Arts, entertainment and recreation                & 154 & 138 & 16.55 & 89.6 \\
Other service activities                          & 154 & 134 & 12.84 & 87.0 \\
Accommodation and food service activities         & 154 & 134 &  9.22 & 87.0 \\
Real estate activities                            & 154 & 129 & 11.47 & 83.8 \\
Education                                         & 154 & 127 &  9.95 & 82.5 \\
Administrative and support service activities     & 154 & 117 &  7.43 & 76.0 \\
Public administration and defence                 & 126 &  86 &  8.07 & 68.3 \\
Mining and quarrying                              & 148 &  96 &  6.94 & 64.9 \\
Electricity, gas, steam and air supply            & 142 &  92 &  8.48 & 64.8 \\
Transportation and storage                        & 154 &  78 &  3.79 & 50.6 \\
Construction                                      & 154 &  55 & $-$2.72 & 35.7 \\
Water supply; sewerage \& waste management        & 152 &  43 & $-$0.28 & 28.3 \\
\bottomrule
\end{tabular}
\end{adjustbox}
\begin{tablenotes}
\footnotesize
\item \textit{Notes:} Exposure\,=\,share of country-year cells with gender pay
  gap $\geq 5\%$ (Directive joint-pay-assessment threshold).  Single-letter
  NACE Rev.~2 sections.  Agriculture (A) excluded (zero usable observations).
  Sample: 27~EU member states, 2019--2024.  Sorted by exposure share
  descending.
\end{tablenotes}
\end{threeparttable}
\end{table}

% ------- TABLE 6c ----------------------------------------------
\begin{table}[htbp]
\caption{Country-Level Compliance Exposure}
\label{tab:countryexposure}
\begin{threeparttable}
\begin{adjustbox}{max width=\textwidth}
\begin{tabular}{llrrrrrr}
\toprule
Code & Country & $N$ usable & $N \geq 5\%$ & Mean GPG (\%) & HPI & ERS & Exposure (\%) \\
\midrule
AT & Austria     & 114 & 114 & 17.70 &  88 &  97 & \textbf{100.0} \\
EE & Estonia     & 124 & 117 & 17.67 &  69 & 100 & 94.4 \\
CZ & Czechia     & 126 & 118 & 16.52 &  93 &  99 & 93.7 \\
LV & Latvia      & 126 & 117 & 17.65 &  76 &  71 & 92.9 \\
EL & Greece      &  57 &  52 & 12.16 &  42 &  74 & 91.2 \\
NL & Netherlands & 126 & 114 & 13.03 & 100 &  56 & 90.5 \\
DE & Germany     & 126 & 113 & 15.24 &  98 &  81 & 89.7 \\
FI & Finland     & 126 & 113 & 12.33 &  44 &  87 & 89.7 \\
IE & Ireland     &  36 &  32 & 11.16 &  79 &  44 & 88.9 \\
DK & Denmark     & 126 & 110 & 10.61 &  55 &  74 & 87.3 \\
CY & Cyprus      & 126 & 106 & 13.96 &  90 &  68 & 84.1 \\
ES & Spain       & 126 & 106 &  8.99 &  48 &  39 & 84.1 \\
SK & Slovakia    & 126 & 104 & 14.01 &  45 &  83 & 82.5 \\
IT & Italy       & 126 & 102 & 13.48 &  71 &  34 & 81.0 \\
HU & Hungary     & 126 & 102 & 14.60 &  88 &  90 & 81.0 \\
FR & France      & 126 &  99 & 10.17 &  67 &  67 & 78.6 \\
LT & Lithuania   & 126 &  98 & 12.46 &  78 &  52 & 77.8 \\
MT & Malta       & 106 &  80 &  9.57 &  80 &  27 & 75.5 \\
HR & Croatia     & 126 &  95 &  9.68 &  68 &  32 & 75.4 \\
BG & Bulgaria    & 126 &  95 &  9.10 &  81 &  63 & 75.4 \\
PL & Poland      & 126 &  95 & 10.62 &  82 &  23 & 75.4 \\
PT & Portugal    & 114 &  84 & 11.00 &  69 &  39 & 73.7 \\
SI & Slovenia    & 126 &  91 &  7.11 &  99 &  42 & 72.2 \\
SE & Sweden      & 126 &  79 &  6.44 &  56 &  58 & 62.7 \\
LU & Luxembourg  & 114 &  63 &  6.61 &  54 &   0 & 55.3 \\
BE & Belgium     & 114 &  56 &  4.87 &  85 &   4 & 49.1 \\
RO & Romania     & 126 &  60 &  4.85 &  42 &  18 & 47.6 \\
\bottomrule
\end{tabular}
\end{adjustbox}
\begin{tablenotes}
\footnotesize
\item \textit{Notes:} Exposure\,=\,share of country-sector-year cells with
  gender pay gap $\geq 5\%$.  $N$ usable counts cells with non-missing gender
  pay gap data across all NACE sectors and years 2019--2024.  HPI and ERS are
  time-invariant latest-period scores.  Sorted by exposure share descending.
\end{tablenotes}
\end{threeparttable}
\end{table}

% ------- FIGURE 5 ----------------------------------------------
\begin{figure}[htbp]
  \centering
  \includegraphics[width=0.88\textwidth]{fig5_compliance_exposure}
  \caption{Compliance Exposure by Sector. Share of country-year cells with
    gender pay gap $\geq 5\%$ (Directive joint-pay-assessment threshold) by
    NACE Rev.~2 section.  Financial and insurance activities and Information \&
    Communication are at 100\% exposure across the full 2019--2024 sample.
    Sample: 27~EU member states.}
  \label{fig:exposure}
\end{figure}

\subsubsection{Out-of-sample validation}
\label{subsec:oos_validation}

The sector-exposure pattern above is descriptive: Table~\ref{tab:sectorexposure}
shows that sector predicts exposure better than tightness \emph{within} the
2019--2024 sample used to construct it.  To test whether this pattern
generalises rather than merely describing the training data, we validate it
out of sample.  We fit three logistic regression classifiers of
$P(\text{GPG}_{st} \geq 5\%)$ on country-sector-year cells from 2019--2022
only ($N=1{,}807$): a full model (sector fixed effects, country fixed
effects, HPI, and employment rate), a tightness-only model (HPI and
employment rate, no sector information), and a sector-only model (sector
fixed effects only, no tightness information).  Each model is then used to
predict exposure on the 2023--2024 cells ($N=917$), which were held out
entirely during fitting.

Table~\ref{tab:oosvalidation}, Panel~A reports the results.  The full model
achieves a test-set AUC of 0.862; the sector-only model, with no tightness
information at all, achieves 0.805 --- capturing most of the full model's
discriminative power.  The tightness-only model achieves an AUC of 0.560,
barely above the 0.500 no-skill baseline.  The accuracy column requires a
caution the AUC does not: the tightness-only model's reported accuracy of
76.1\% is identical to the test-set base rate (698 of 917 held-out cells,
or 76.1\%, are actually exposed), and its confusion matrix confirms why ---
it predicts ``exposed'' for every single test observation, with zero true
negatives and zero false negatives.  This is a degenerate classifier with no
discriminative content, not a moderately accurate one; AUC, not accuracy, is
the informative statistic for this model.

Panel~B confirms the AUC gaps are not an artefact of one train/test split.
Both a DeLong test and a 2,000-iteration bootstrap agree closely: the
full-model advantage over tightness-only ($\Delta\text{AUC}=0.302$, 95\% CI
$[0.256, 0.348]$) and the sector-only advantage over tightness-only
($\Delta\text{AUC}=0.245$, 95\% CI $[0.189, 0.301]$) are both significant at
$p<0.0001$.  This converts the descriptive claim in
Section~\ref{subsec:exposure} into a validated predictive one: sector
membership, not labour-market tightness, generalises to compliance-threshold
exposure in years the model never saw during training.

% ------- TABLE: OUT-OF-SAMPLE VALIDATION ------------------------
\begin{table}[htbp]
\caption{Out-of-Sample Validation of Sector-Based Exposure Mapping}
\label{tab:oosvalidation}
\begin{threeparttable}
\smallskip
\textit{Panel A: Test-set performance (trained on 2019--2022, tested on 2023--2024)}
\smallskip
\begin{tabular}{lccc}
\toprule
Model & $N_{\text{train}}$ & $N_{\text{test}}$ & AUC \quad Accuracy (\%) \\
\midrule
Full (sector + country FE + HPI + employment rate) & 1,807 & 917 & 0.862 \quad 83.9 \\
Sector-only (sector fixed effects, no tightness)    & 1,807 & 917 & 0.805 \quad 80.7 \\
Tightness-only (HPI + employment rate, no sector)   & 1,807 & 917 & 0.560 \quad 76.1\sym{a} \\
\bottomrule
\end{tabular}

\bigskip
\textit{Panel B: Significance of AUC gaps vs.\ tightness-only}
\smallskip
\begin{tabular}{llccc}
\toprule
Comparison & Method & $\Delta$AUC & 95\% CI & $p$-value \\
\midrule
Full vs.\ tightness-only         & DeLong    & 0.302 & [0.256, 0.348] & $<$0.001 \\
Full vs.\ tightness-only         & Bootstrap & 0.303 & [0.258, 0.348] & $<$0.001 \\
Sector-only vs.\ tightness-only  & DeLong    & 0.245 & [0.189, 0.301] & $<$0.001 \\
Sector-only vs.\ tightness-only  & Bootstrap & 0.246 & [0.192, 0.302] & $<$0.001 \\
\bottomrule
\end{tabular}
\begin{tablenotes}
\footnotesize
\item \textit{Notes:} Training set: 1,807 country-sector-year cells, 2019--2022.
  Test set: 917 held-out cells, 2023--2024, never used in fitting.  Test-set
  base rate: 698/917 cells (76.1\%) are actually $\geq 5\%$ exposed.
  Bootstrap: 2,000 resamples of the test set.
  \sym{a} The tightness-only model's accuracy is identical to the test-set
  base rate because it predicts ``exposed'' for every test observation
  (0 true negatives, 0 false negatives); its accuracy figure reflects no
  discriminative skill, unlike its AUC.
\end{tablenotes}
\end{threeparttable}
\end{table}

\subsection{Combined risk quadrant}
\label{subsec:quadrant}

Figure~\ref{fig:quadrant} plots each country's Hiring Pressure Index against its
Equity Risk Score, divided into four quadrants at the midpoint of each axis.
This classification is descriptive triage, not a causal or predictive claim:
given Sections~\ref{subsec:pooled_fe}--\ref{subsec:mechanism}'s null tightness
result, a country's position in this quadrant should not be read as implying that
its labour-market tightness is driving its equity risk level.  Germany, Czechia,
Hungary, and Austria fall in the ``Priority Intervention'' quadrant (high
tightness, high equity risk); the addition of Austria to this group reflects a
data-completeness correction rather than any change in the tightness--pay-gap
relationship itself.  The quadrant is presented as a practical starting point for
prioritising compliance review, to be read alongside --- not instead of --- the
sector-level exposure evidence in Section~\ref{subsec:exposure}.
Figure~\ref{fig:prioritycases} provides a closer look at the priority-quadrant
countries.

% ------- FIGURE 2 ----------------------------------------------
\begin{figure}[htbp]
  \centering
  \includegraphics[width=0.82\textwidth]{fig2_combined_risk_quadrant}
  \caption{Combined Risk Quadrant. Each country is plotted by Hiring Pressure
    Index (x-axis, higher = tighter labour market) against Equity Risk Score
    (y-axis, higher = larger gender pay gap).  Quadrant boundaries at the
    midpoint of each axis.  Countries in the top-right quadrant (high tightness,
    high equity risk) are the ``Priority Intervention'' group: Germany, Czechia,
    Hungary, and Austria.  \textit{Note:} Tightness is not a reliable
    within-country predictor of equity risk (Section~\ref{subsec:pooled_fe});
    the quadrant is a descriptive triage tool only.}
  \label{fig:quadrant}
\end{figure}

% ------- FIGURE 4 ----------------------------------------------
\begin{figure}[htbp]
  \centering
  \includegraphics[width=0.82\textwidth]{fig4_priority_cases}
  \caption{Priority Intervention Cases. Hiring Pressure Index and Equity Risk
    Score for countries in the top-right quadrant of Figure~\ref{fig:quadrant}
    (Germany, Czechia, Hungary, Austria). Both indices are time-invariant
    latest-period scores; see Section~\ref{subsec:composites}.}
  \label{fig:prioritycases}
\end{figure}

% ============================================================
%  4. DISCUSSION
% ============================================================
\section{Discussion}
\label{sec:discussion}

The central empirical fact of this paper is a contrast, not a single number: a
labour-market-tightness/gender-pay-gap correlation that looks real and
policy-relevant when tested only across countries evaporates when tested within
countries over time.  This pattern is consistent with the correlation being
driven by time-invariant, country-level characteristics --- industrial
composition, wage-setting institutions, historical labour-market structure ---
that happen to be correlated with both a country's typical tightness level and
its typical pay gap level, without either causing the other in the way
competitive equalisation theory would require.  We tested two specific,
well-motivated candidates for such a confound.  The gender employment
participation gap, following \citet{OlivettiPetrongolo2008}, does not explain
the pooled association (Section~\ref{subsec:mechanism}).  Collective bargaining
coverage, motivated by \citet{VaccaroEtAl2024}'s finding that pay-transparency
legislation is effective only where worker representation exists, does not
moderate the within-country null on our pre-specified test, though a secondary,
more exploratory cut suggests the relationship may differ by coverage level and
warrants further investigation with richer data (Section~\ref{subsec:bargaining}).
We additionally cross-checked the central null against a source of variation
with no obvious channel back to national gender-discrimination institutions ---
the differential severity of the 2020 COVID-19 lockdown shock across countries
--- and found the same null result (Section~\ref{subsec:covid_event}), with the
caveat that this design cannot rule out pre-existing divergent trends given only
one pre-treatment year in the panel.  We cannot rule out other unobserved
confounders we did not test directly, such as industrial composition or
childcare policy regimes, and this remains a genuine limitation rather than a
settled matter.

The policy implication is direct and, we think, more useful than a confirmed
causal finding would have been.  The Directive is being transposed on the
assumption that market forces alone are not closing gender pay gaps; nothing in
this paper contradicts that starting premise.  What this paper does show is that
labour market tightness --- whether measured by employment rate or job vacancy
rate --- is not a reliable signal, in either direction, for predicting where
compliance risk under the Directive will concentrate.  A regulator or employer
using ``how tight is this country's labour market'' as a triage heuristic for
compliance risk would be using a heuristic with essentially no predictive power
in this data (Section~\ref{subsec:exposure}).  This is not merely an in-sample
pattern: a classifier trained only on 2019--2022 data and tested on 2023--2024
confirms that sector-based prediction substantially and significantly
outperforms tightness-based prediction out of sample
(Section~\ref{subsec:oos_validation}).  Sector, by contrast, is a strong
and stable predictor: Finance and Information and Communication show 100\%
exposure to the Directive's 5\% threshold across the entire 2019--2024 sample, a
pattern that holds regardless of country-level tightness conditions.

A further limitation concerns the temporal alignment between this analysis and
the Directive's implementation timeline.  Employer-level pay-gap reporting under
the Directive begins in 2027 for the largest firms (250 or more employees),
reporting on the 2026 calendar year.  This paper's analysis uses Eurostat's
aggregate, macro-level gender pay gap series, which will continue to lag
firm-level disclosures once those begin to accumulate.  The question of whether
firm-level Directive disclosures, once available, show the same sector
concentration and the same absence of a tightness relationship found here at the
macro level is a natural and, we think, valuable extension once 2027 reporting
data exists; it is a distinct empirical question from the one this paper answers,
and we do not attempt to anticipate its answer here.

Finally, we note explicitly that the principal components validation of the
Hiring Pressure Index's weighting scheme (Appendix~\ref{sec:appendix}) is weak
on the full, corrected panel (Spearman $\rho=0.657$, $p=0.156$) and does not
meet conventional significance thresholds.  We report this transparently rather
than omitting it or overstating it, consistent with this paper's broader stance
that a claim which looked stronger on a partial or earlier sample should be
re-tested, not carried forward uncritically, once fuller data is available.  This
diagnostic result does not affect the paper's three main contributions, none of
which depend on the PCA validation holding.

% ============================================================
%  5. CONCLUSION
% ============================================================
\section{Conclusion}
\label{sec:conclusion}

This paper set out to test, using panel methods appropriate to the question, an
empirical premise implicit in the EU Pay Transparency Directive: that labour
market tightness bears a reliable relationship to the gender pay gap.  Using a
27-country Eurostat panel spanning 2019--2024, we find that a positive,
statistically significant cross-sectional association between employment rate and
the gender pay gap does not survive the introduction of country and year fixed
effects, and remains null across every robustness check and mechanism control we
apply.  Labour market tightness, measured either by employment rate or job
vacancy rate, is not a reliable predictor of gender pay gap levels within a given
country over time, nor is it a useful predictor of which country-sector cells
will exceed the Directive's 5\% compliance threshold.

What is robust in this data is sector.  Gender pay gaps differ sharply and
consistently across NACE sectors, with Financial and insurance activities and
Information and Communication activities showing persistently high gaps and 100\%
exposure to the Directive's compliance threshold across the entire sample period,
regardless of the tightness conditions in the country where they operate.  We
propose sector-based exposure mapping, rather than any labour-market-tightness
heuristic, as the more actionable basis for compliance prioritisation as employers
and regulators prepare for the Directive's 2027 reporting requirements.

% ============================================================
%  APPENDIX A
% ============================================================
\appendix
\section{PCA Validation of the Hiring Pressure Index (Diagnostic)}
\label{sec:appendix}

The Hiring Pressure Index combines six underlying Eurostat signals --- job
vacancy rate, unemployment rate, employment rate, labour market slack, and two
labour-flow measures --- using hand-set weights (Section~\ref{subsec:composites}).
To assess whether these weights track the dominant axis of variation in the
underlying data, we run a principal components analysis on the six standardised
input signals across 172 country-year observations spanning 25 countries with
complete data on all six signals.

The first principal component explains 46.4\% of total variance (eigenvalue\,=\,2.80),
with the dominant loading on unemployment rate; the second component explains a
further 24.5\% (cumulative 70.9\%; Table~\ref{tab:pca_eigen}).  Comparing the
empirical HPI formula weights, normalised to $[-1, 1]$, against the first
component's loadings (Table~\ref{tab:pca_loadings}), with component sign anchored
so that higher values correspond to higher vacancy rate, the Spearman rank
correlation is $\rho=0.657$ ($p=0.156$, $N=6$ signals).  This does not meet
conventional significance thresholds.

We report this result as a transparent diagnostic rather than as a validated
methodological contribution.  An earlier analysis on a smaller, 22-country sample
(since corrected for a sector-classification data issue affecting five countries)
found a stronger and nominally significant correlation ($\rho=0.886$, $p=0.019$).
The weaker result on the full, corrected 27-country sample is the one we report
and rely on; we do not carry forward the earlier, more favourable estimate.
This appendix result does not affect any of the paper's three main contributions
(Section~\ref{sec:intro}), each of which is independent of whether the HPI
weighting scheme is formally PCA-validated.

% ------- TABLE A1 ----------------------------------------------
\begin{table}[htbp]
\caption{PCA Results: Eigenvalues and Variance Explained}
\label{tab:pca_eigen}
\begin{threeparttable}
\begin{tabular}{lrrrr}
\toprule
Component & Eigenvalue & Var.\ explained (\%) & Cumulative (\%) \\
\midrule
PC1 & 2.802 & 46.44 & 46.44 \\
PC2 & 1.479 & 24.51 & 70.94 \\
PC3 & 0.939 & 15.56 & 86.50 \\
PC4 & 0.547 &  9.06 & 95.56 \\
PC5 & 0.200 &  3.31 & 98.87 \\
PC6 & 0.068 &  1.13 & 100.00 \\
\bottomrule
\end{tabular}
\begin{tablenotes}
\footnotesize
\item \textit{Notes:} PCA on six standardised HPI input signals
  (job vacancy rate, unemployment rate, employment rate, labour slack rate,
  flow-to-employment, flow-to-inactivity) across 172 country-year observations,
  25 countries with complete data on all six signals.
\end{tablenotes}
\end{threeparttable}
\end{table}

% ------- TABLE A2 ----------------------------------------------
\begin{table}[htbp]
\caption{PCA Factor Loadings on First Four Principal Components}
\label{tab:pca_loadings}
\begin{threeparttable}
\begin{tabular}{lrrrr}
\toprule
Signal & PC1 & PC2 & PC3 & PC4 \\
\midrule
Job vacancy rate          &  0.399 &  0.338 & $-$0.093 &  0.830 \\
Unemployment rate         & $-$0.519 &  0.227 & $-$0.367 &  0.064 \\
Employment rate           &  0.506 & $-$0.001 & $-$0.326 & $-$0.418 \\
Labour slack rate         & $-$0.446 &  0.493 & $-$0.227 & $-$0.050 \\
Flow to employment        &  0.340 &  0.558 & $-$0.324 & $-$0.306 \\
Flow to inactivity        &  0.027 &  0.529 &  0.770 & $-$0.189 \\
\bottomrule
\end{tabular}
\begin{tablenotes}
\footnotesize
\item \textit{Notes:} Loadings on first four principal components from PCA on
  six standardised HPI input signals.  PC1 eigenvalue\,=\,2.802 (46.4\% of
  variance); PC2 eigenvalue\,=\,1.479 (24.5\%).
  Spearman rank correlation between HPI formula weights (normalised) and PC1
  loadings: $\rho=0.657$, $p=0.156$, $N=6$ signals (not significant).
\end{tablenotes}
\end{threeparttable}
\end{table}

% ------- FIGURE A1 + A2 ---------------------------------------
\begin{figure}[htbp]
  \centering
  \includegraphics[width=0.75\textwidth]{fig_a1_scree_plot}
  \caption{Scree Plot: Principal Components of HPI Input Signals. Eigenvalues
    for each of the six principal components from PCA on the standardised HPI
    inputs (172 country-year observations). The elbow suggests two dominant
    components (PC1: 46.4\%, PC2: 24.5\%; cumulative 70.9\%).}
  \label{fig:scree}
\end{figure}

\begin{figure}[htbp]
  \centering
  \includegraphics[width=0.80\textwidth]{fig_a2_pca_biplot}
  \caption{PCA Biplot: Six HPI Input Signals. Biplot of the first two principal
    components across 172 country-year observations (25 countries with complete
    data). Vectors show the loadings of each input signal on PC1 and PC2.
    Spearman $\rho=0.657$ ($p=0.156$) between HPI formula weights and PC1
    loadings; the HPI weighting is not formally PCA-validated at conventional
    significance thresholds.}
  \label{fig:biplot}
\end{figure}

\subsection*{Robustness to alternative weighting schemes}
\label{subsec:hpi_weight_robustness}

The comparison above tests whether the HPI formula's \emph{weights} track the
dominant axis of variation among the six input signals, and finds only a weak,
non-significant correlation.  This is a narrower test than asking whether the
\emph{resulting index scores} are themselves sensitive to the specific
weighting scheme chosen.  We test this separately by constructing three
alternative country-year HPI variants from the same six underlying signals ---
an equal-weighting scheme, a scheme that heavily weights job vacancy rate
relative to the other five signals, and a scheme using the first principal
component's own loadings as weights --- and computing each variant's Spearman
rank correlation with the HPI actually used throughout this paper, across the
full 189 country-year panel.

All three alternative weighting schemes correlate strongly and significantly
with the reported HPI: $\rho=0.734$ ($p<0.001$) for equal weighting,
$\rho=0.889$ ($p<0.001$) for the vacancy-heavy scheme, and $\rho=0.838$
($p<0.001$) for the PCA-loadings scheme (Table~\ref{tab:hpiweightrobustness}).
This is a materially different result from the weight-vector comparison above,
and a reassuring one: even though the paper's hand-set weights do not closely
track the first principal component when compared as abstract weight vectors,
the country-year rankings the HPI actually produces are not sensitive to that
choice --- three plausible alternative weighting schemes, including one built
directly from PC1, all reproduce substantially the same ordering.  The
substantive results in this paper that depend on HPI (Sections~\ref{subsec:exposure}
and~\ref{subsec:quadrant}) are therefore unlikely to be an artefact of the
specific hand-set weights in Section~\ref{subsec:composites}.

% ------- TABLE A3 -----------------------------------------------
\begin{table}[htbp]
\caption{HPI Robustness to Alternative Weighting Schemes}
\label{tab:hpiweightrobustness}
\begin{threeparttable}
\begin{tabular}{lccc}
\toprule
HPI variant & Spearman $\rho$ & $p$-value & $N$ \\
\midrule
Equal weighting            & 0.734 & $<$0.001 & 189 \\
Vacancy-rate-heavy weighting & 0.889 & $<$0.001 & 189 \\
PCA-loadings weighting      & 0.838 & $<$0.001 & 189 \\
\bottomrule
\end{tabular}
\begin{tablenotes}
\footnotesize
\item \textit{Notes:} Each variant is constructed from the same six
  standardised input signals as the reported HPI (Section~\ref{subsec:composites})
  under an alternative weighting scheme, then rank-correlated (Spearman) against
  the HPI actually used throughout this paper, across the full 189 country-year
  panel.  This tests sensitivity of the resulting index scores to the weighting
  scheme, distinct from the weight-vector-level comparison against PC1 loadings
  above ($N=6$ signals).
\end{tablenotes}
\end{threeparttable}
\end{table}

% ============================================================
%  BIBLIOGRAPHY
% ============================================================
\newpage
\bibliographystyle{apalike}
\bibliography{references}

% ============================================================
%  END MATTER
% ============================================================
\newpage
% ============================================================
%  APPENDIX B
% ============================================================
\section{A Pre-Registered Design for Future Data: Exploiting Staggered
         Transposition Timing}
\label{sec:appendixB}

The Directive's 7~June~2026 transposition deadline was met unevenly across
the EU27: some member states transposed on schedule, others delayed
implementation to 2027 or later, and several remained unresolved at the
deadline.\footnote{Rolling summaries of national transposition status are
maintained by international law firms tracking Directive compliance (e.g.,
Baker McKenzie, Littler, and PwC publish periodically updated implementation
trackers); we do not rely on any single tracker as authoritative and note
that transposition status continues to evolve after this paper's writing.}
This heterogeneity in \emph{when} the Directive's compliance obligations
become legally binding is a genuine source of quasi-experimental variation
distinct from anything used in Sections~\ref{sec:results}: rather than
asking whether tightness predicts the gender pay gap, it asks whether
anticipation of binding compliance obligations changes employer or
labour-market behaviour around each country's actual transposition date.

We cannot test this design in the present paper.  Eurostat's gender pay gap
series (\texttt{earn\_gr\_gpgr2}) is restricted to reference years through
2024 (Section~\ref{subsec:data_sources}), and firm-level disclosures under
the Directive do not begin until 2027, reporting on the 2026 calendar year.
There is therefore no post-transposition outcome data available for any
country's transposition date as of this writing.  We record the design here
so that it can be executed once the relevant data exists, rather than leaving
it undocumented, and to make clear that the absence of a result in this
appendix reflects a data-availability constraint rather than an oversight.

\subsection*{Proposed specification}

Let $T_i$ denote country $i$'s actual transposition date (or, where formal
transposition remains unresolved, the date national courts began applying
the Directive's substantive standard to existing discrimination claims,
whichever is earlier and better documented).  A staggered-adoption
event-study specification, following the standard structure used elsewhere
in this paper, would take the form

\begin{equation}
Y_{it} = \sum_{k \neq -1} \theta_k \,\mathbb{1}[t - T_i = k]
  + \gamma_i + \delta_t + \varepsilon_{it},
\label{eq:transposition}
\end{equation}

\noindent where $Y_{it}$ is the outcome of interest (candidates include the
gender pay gap itself once post-2024 data exists, or a firm-behaviour proxy
such as job-posting salary-range disclosure rates once available), $k$
indexes event time relative to each country's own transposition date, and
$k=-1$ is the omitted reference period.  Given the well-documented
econometric concerns with staggered-adoption designs under
treatment-effect heterogeneity, any execution of Equation~\ref{eq:transposition}
should use an estimator robust to this issue (e.g., \citealt{CallawaySantAnna2021}
or \citealt{SunAbraham2021}) rather than standard two-way fixed effects, and
should assess whether the timing of transposition is itself correlated with
each country's pre-existing tightness or gender-pay-gap trajectory before
treating $T_i$ as exogenous.

We flag this design explicitly as an intended extension rather than a
finding, consistent with this paper's treatment of other results that are
not yet fully supported by available data (Section~\ref{subsec:bargaining}).

% ============================================================
%  APPENDIX C
% ============================================================
\section{Full Pairwise Comparisons: Finance Versus All Other Sectors}
\label{sec:appendixC}

Section~\ref{subsec:sector} reports that every pairwise comparison of
Financial and insurance activities against every other NACE section remains
significant after Bonferroni correction for multiple comparisons, but does not
tabulate the underlying comparisons.  Table~\ref{tab:financepairwise} reports
all 17 pairwise mean-difference tests in full.  Every comparison is significant
at $p<0.001$ after Bonferroni adjustment, with mean differences ranging from
27.7 percentage points (versus Construction) to 5.5 percentage points (versus
Information and Communication, itself the second-highest-gap sector in the
sample).

% ------- TABLE C1 -----------------------------------------------
\begin{table}[htbp]
\caption{Pairwise Comparisons: Finance Versus All Other Sectors (Bonferroni-Adjusted)}
\label{tab:financepairwise}
\begin{threeparttable}
\begin{adjustbox}{max width=\textwidth}
\begin{tabular}{lrrrr}
\toprule
Sector (vs.\ Finance) & Mean difference (pp) & $t$-statistic & $p$-value & $p$-value (Bonferroni) \\
\midrule
Construction                                        & 27.67 & 25.87 & $<$0.001 & $<$0.001 \\
Water supply; sewerage \& waste management          & 25.23 & 22.34 & $<$0.001 & $<$0.001 \\
Transportation and storage                          & 21.17 & 24.08 & $<$0.001 & $<$0.001 \\
Mining and quarrying                                & 18.01 & 17.78 & $<$0.001 & $<$0.001 \\
Administrative and support service activities       & 17.52 & 18.15 & $<$0.001 & $<$0.001 \\
Public administration and defence                   & 16.88 & 23.76 & $<$0.001 & $<$0.001 \\
Electricity, gas, steam and air supply               & 16.48 & 20.09 & $<$0.001 & $<$0.001 \\
Accommodation and food service activities           & 15.74 & 24.77 & $<$0.001 & $<$0.001 \\
Education                                           & 15.01 & 21.29 & $<$0.001 & $<$0.001 \\
Real estate activities                              & 13.49 & 15.99 & $<$0.001 & $<$0.001 \\
Other service activities                            & 12.12 & 14.12 & $<$0.001 & $<$0.001 \\
Professional, scientific and technical activities   &  8.65 & 11.04 & $<$0.001 & $<$0.001 \\
Arts, entertainment and recreation                  &  8.41 &  6.91 & $<$0.001 & $<$0.001 \\
Manufacturing                                       &  8.21 & 10.51 & $<$0.001 & $<$0.001 \\
Human health and social work activities             &  7.89 &  9.04 & $<$0.001 & $<$0.001 \\
Wholesale and retail trade                          &  6.72 &  9.53 & $<$0.001 & $<$0.001 \\
Information and communication                       &  5.49 &  7.26 & $<$0.001 & $<$0.001 \\
\bottomrule
\end{tabular}
\end{adjustbox}
\begin{tablenotes}
\footnotesize
\item \textit{Notes:} Mean difference\,=\,Finance's mean gender pay gap minus
  the comparison sector's mean gender pay gap, in percentage points.  Two-sample
  $t$-tests; $p$-values Bonferroni-adjusted for 17 simultaneous comparisons.
  Sample: 27~EU member states, 2019--2024 (single-letter NACE Rev.~2 sections;
  Agriculture excluded, zero usable observations).
\end{tablenotes}
\end{threeparttable}
\end{table}

\section*{Data Availability Statement}

All data underlying this study are publicly available from Eurostat's JSON-stat
API (\url{https://ec.europa.eu/eurostat/web/json-and-unicode-web-services}).  The
exact datasets, pull scripts, and versioned Parquet extracts used in this study,
along with all analysis code and generated tables and figures, are archived at
\url{https://github.com/SVamseekar/workforceguardai} and at Zenodo
(DOI:\,\href{https://doi.org/10.5281/zenodo.20455974}{10.5281/zenodo.20455974};
to be updated with a replication-package-specific record prior to submission).

\section*{Funding}

This research received no external funding.

\section*{Conflicts of Interest}

The author declares no conflicts of interest.

\section*{Author Note}

Marti Soura Vamseekar is an independent researcher.  M.Sc.\ Data Science,
University of Greenwich, London, UK\@.  Correspondence:
\href{mailto:martisoura@gmail.com}{martisoura@gmail.com}.
ORCID:\;\href{https://orcid.org/0009-0005-4884-1292}{0009-0005-4884-1292}.

\end{document}
